% 平行性质：a∥b ⇒ 三种角关系
\documentclass[tikz,border=4pt]{standalone}
\usepackage{ctex}
\usepackage{amsmath}
\usetikzlibrary{calc, arrows.meta, decorations.markings}
\begin{document}
\begin{tikzpicture}[scale=0.85, thick,
  arrowone/.style={postaction={decorate, decoration={markings,
    mark=at position 0.5 with {\arrow{Stealth}}}}},
  arrowtwo/.style={postaction={decorate, decoration={markings,
    mark=at position 0.45 with {\arrow{Stealth}};
    mark=at position 0.55 with {\arrow{Stealth}}}}}]
  \foreach \shift/\caption/\labelOne/\labelTwo in {
    {0}/{同位角相等}/{P-下}/{Q-下},
    {5.5}/{内错角相等}/{P-下}/{Q-上},
    {11}/{同旁内角互补}/{P-下}/{Q-上同侧}
  } {
    \begin{scope}[shift={(\shift,0)}]
      \draw[arrowone] (-1.5,1.8) -- (2.5,1.8);
      \draw[arrowtwo] (-1.5,0) -- (2.5,0);
      \draw (-0.5,2.6) -- (1.1,-0.8);
    \end{scope}
  }
  % 同位
  \begin{scope}[shift={(0,0)}]
    \coordinate (P) at (0,1.8); \coordinate (Q) at (0.6,0);
    \draw[red] ($(P)+(0.4,0)$) arc (0:-65:0.4);
    \draw[red] ($(Q)+(0.4,0)$) arc (0:-65:0.4);
    \node[below, font=\small] at (0.5,-1.2) {$a\parallel b\Rightarrow$ 同位角相等};
  \end{scope}
  % 内错
  \begin{scope}[shift={(5.5,0)}]
    \coordinate (P) at (0,1.8); \coordinate (Q) at (0.6,0);
    \draw[red] ($(P)+(0.4,0)$) arc (0:-65:0.4);
    \draw[red] ($(Q)+(-0.4,0)$) arc (180:115:0.4);
    \node[below, font=\small] at (0.5,-1.2) {$a\parallel b\Rightarrow$ 内错角相等};
  \end{scope}
  % 同旁内
  \begin{scope}[shift={(11,0)}]
    \coordinate (P) at (0,1.8); \coordinate (Q) at (0.6,0);
    \draw[red] ($(P)+(0.4,0)$) arc (0:-65:0.4);
    \draw[red] ($(Q)+(0.4,0)$) arc (0:115:0.4);
    \node[below, font=\small] at (0.5,-1.2) {$a\parallel b\Rightarrow$ 同旁内角互补};
  \end{scope}
\end{tikzpicture}
\end{document}
